\documentclass[conference]{IEEEtran}
\usepackage{cite}
\usepackage{amsmath,amssymb,amsfonts}
\usepackage{algorithmic}
\usepackage{graphicx}
\usepackage{textcomp}
\usepackage{xcolor}
\usepackage{booktabs}
\usepackage{float}
\usepackage{hyperref}

\def\BibTeX{{\rm B\kern-.05em{\sc i\kern-.025em b}\kern-.08em
    T\kern-.1667em\lower.7ex\hbox{E}\kern-.125emX}}

\begin{document}

\title{A Geo AI - Based Framework For Geospatial Flood Risk Mapping And Short-Term Rainfall Prediction For Urban Waterlogging Prevention}

\author{
    \begin{minipage}{0.45\textwidth}
        \centering
        \textbf{Jacob Mathew} \\
        \textit{Department of Artificial Intelligence} \\
        \textit{Providence College of Engineering} \\
        Chengannur, Kerala, India \\
        jacobmathew627@gmail.com
    \end{minipage}
    \hfill
    \begin{minipage}{0.45\textwidth}
        \centering
        \textbf{Gowri Shankar P} \\
        \textit{Department of Artificial Intelligence} \\
        \textit{Providence College of Engineering} \\
        Chengannur, Kerala, India \\
        gowrishankars706@gmail.com
    \end{minipage}
    \vspace{0.7cm} \\
    \begin{minipage}{0.45\textwidth}
        \centering
        \textbf{Siddharth Biju} \\
        \textit{Department of Artificial Intelligence} \\
        \textit{Providence College of Engineering} \\
        Chengannur, Kerala, India \\
        siddharthbiju18@gmail.com
    \end{minipage}
    \hfill
    \begin{minipage}{0.45\textwidth}
        \centering
        \textbf{Anoop P.P.} \\
        \textit{Assistant Professor, Dept. of AI} \\
        \textit{Providence College of Engineering} \\
        Chengannur, Kerala, India \\
        anoop.p@providence.edu.in
    \end{minipage}
}

\maketitle

\begin{abstract}
Urban waterlogging during intense monsoon events causes widespread socioeconomic disruption in densely populated coastal cities. Ernakulam, Kerala, India, has been severely impacted by recurrent flooding, notably during the 2018 disaster. While hydrodynamic models provide high accuracy, they require comprehensive, proprietary underground drainage blueprints that are frequently unavailable in rapidly developing regions. This paper presents a GeoAI framework designed to bypass these data constraints through a two-stage predictive pipeline. The first stage, Short-Term Rainfall Prediction, purposefully circumvents the immense computational overhead of training redundant, custom meteorological models by directly integrating forecasted telemetry from established, state-of-the-art third-party Weather APIs. These discrete API telemetry points are spatially interpolated via Inverse Distance Weighting to model impending precipitation loads over a 24-hour horizon. The second stage integrates this dynamic atmospheric load alongside static geographical features---digital elevation model (DEM), slope, topographic wetness index (TWI), surrogate drainage proximity, drainage density, and built-up density---into a 7-channel spatial tensor. An Attention U-Net convolutional neural network is trained on spatially disjoint patches derived from Sentinel-1 SAR inundation labels to preclude spatial autocorrelation leakage. Validated via spatial block cross-validation and evaluated on an independent 2019 temporal event, the model achieves a competitive F1-score of 0.670 and ROC-AUC of 0.919, outperforming spatially agnostic baselines (Random Forest) and state-of-the-art semantic segmentation networks (DeepLabV3+). Ablation experiments indicate that the API-driven interpolated precipitation and surrogate drainage proximity contribute critically to spatial performance. Deployed via a FastAPI backend, the framework enables continuous real-time flood susceptibility forecasting for localities lacking municipal hydrological data.
\end{abstract}

\begin{IEEEkeywords}
GeoAI, flood susceptibility mapping, Attention U-Net, surrogate drainage, spatial cross-validation, Synthetic Aperture Radar (SAR), deep learning.
\end{IEEEkeywords}

\section{Introduction}

\subsection{Urban Flooding as a Socioeconomic Challenge}
Flooding constitutes one of the most destructive natural hazards globally, amplified by climate change-induced extreme weather events and rapid urbanization \cite{1}. In urban environments, the progressive replacement of permeable surfaces with impervious infrastructure accelerates surface runoff, overwhelming natural hydrological cycles and resulting in severe pluvial flooding \cite{2}. Ernakulam district, located in Kerala, India, represents a developing metropolitan region highly vulnerable to such events. The district experienced devastating inundation during the 2018 southwest monsoon \cite{3}, highlighting the necessity for robust, scalable early warning systems.

\subsection{Limitations of Existing Analytical Methods}
Traditional hydrodynamic models, such as HEC-RAS and SWMM, represent fluid dynamics with high fidelity. However, they require comprehensive underground drainage blueprints, precisely calibrated hydraulic parameters, and immense computational resources. This renders them impractical for many rapidly developing cities where such municipal data is either physically unavailable, outdated, or computationally prohibitive to simulate in real-time during an active crisis \cite{4}.

Statistical and machine learning approaches, including Support Vector Machines (SVM) and Random Forest (RF), have been widely applied due to their scalability. However, these models process pixels independently, discarding the crucial spatial structures---such as continuous valley networks and linear flow paths---that physically dictate flood propagation \cite{5,6}. Consequently, Convolutional Neural Networks (CNNs) have emerged to capture this lost spatial context \cite{7,8}. Yet, many existing CNN flood models treat historical inundation as a static target, lacking parametrizable meteorological inputs required to forecast responses to varying precipitation loads \cite{9,12}.

\subsection{Integrating Short-Term Rainfall Prediction via Weather APIs}
A key methodological integration of this framework is the formal incorporation of Short-Term Rainfall Prediction ($r(t+24)$) directly into the spatial modeling process. Developing custom meteorological forecasting models (such as numerical weather prediction algorithms) requires vast arrays of supercomputing clusters and continuous atmospheric telemetry spanning continental scales \cite{19}, resources far beyond the scope of localized disaster administration. To solve this, the proposed architecture intentionally utilizes existing, highly-optimized third-party Weather APIs (e.g., OpenWeatherMap, IMD data streams) as deterministic inputs. By retrieving this external meteorological telemetry, interpolating it spatially to account for localized variability, and injecting it as a modulating input layer, the system leverages established atmospheric science rather than attempting to redundantly simulate weather phenomena from scratch. This provides an actionable, efficient predictive forecasting engine.

\subsection{Identified Research Gaps}
We identify three distinct research gaps in the current literature, summarized in Table \ref{tab:gaps}. First, traditional high-performing ensemble models lack spatial context. Second, deterministic spatial models frequently rely on official drainage maps that developing nations lack. Third, pure segmentation models often fail to incorporate active, API-driven short-term rainfall predictions into their spatial tensors, limiting them to static hindcasting.

\begin{table}[htbp]
\caption{Comparative Analysis of Recent Flood Mapping Architectures}
\begin{center}
\resizebox{\linewidth}{!}{
\begin{tabular}{|l|c|c|c|c|}
\hline
\textbf{Reference / Method} & \textbf{Spatial} & \textbf{Surrogate} & \textbf{API Rain.} & \textbf{Deployment} \\
\hline
Tehrany et al. \cite{5} & No & No & No & No \\
Konapala et al. \cite{7} & Yes & No & No & No \\
Jiang et al. \cite{17} & Yes & No & No & No \\
B\"{o}hm et al. \cite{12} & Yes & No & No & No \\
\textbf{Proposed Framework} & \textbf{Yes} & \textbf{Yes} & \textbf{Yes} & \textbf{Yes} \\
\hline
\end{tabular}
}
\label{tab:gaps}
\end{center}
\end{table}

\subsection{Contributions}
To address these gaps, the specific contributions of this paper are:
\begin{enumerate}
\item \textbf{Dynamic API Rainfall Integration:} A 7-channel tensor structure where $r(t+24)$ forecasted precipitation, retrieved from commercial Weather APIs and spatially interpolated via Inverse Distance Weighting, links established atmospheric prediction to ground inundation modeling without redundant model training.
\item \textbf{Surrogate Drainage Modeling:} Utilization of a synthetic drainage network derived via D8 flow routing directly from the DEM, eliminating the dependency on proprietary drainage blueprints.
\item \textbf{Rigorous Spatial Validation:} Implementation of spatially disjoint block cross-validation and temporally independent testing to prevent data leakage and prove true predictive generalization.
\item \textbf{End-to-End Operational Deployment:} A complete pipeline featuring an optimized inference backend and interactive dashboard for real-time risk assessment.
\end{enumerate}

\section{Related Work}

\subsection{Machine Learning in Flood Susceptibility}
Early flood susceptibility mapping utilized Multi-Criteria Decision Analysis (MCDA) and bivariate statistics. With the advent of machine learning, ensemble methods like Random Forest gained prominence for pixel-wise mapping \cite{5,6}. These models achieve high ROC-AUC classifications by capturing nonlinear interactions between topology and land cover. However, pixel-wise models are inherently spatially agnostic, often producing fragmented hazard maps and struggling to maintain continuity along linear drainage infrastructure.

\subsection{Deep Learning and Spatiotemporal Advances}
To address these limitations, researchers adopted CNNs. Konapala et al. \cite{7} and Kabir et al. \cite{8} demonstrated that CNNs and U-Nets can recognize the geometric shapes of flood plains, maintaining spatial resolution while extracting deep features. More recently, the field has integrated advanced architectures such as DeepLabV3+ \cite{18} and Swin-Unet to refine boundary delineation through atrous spatial pyramid pooling and attention mechanisms \cite{17}. Additionally, spatiotemporal frameworks utilizing ConvLSTM or Graph Neural Networks (GNNs) have been proposed to model flood wave propagation over time \cite{19,20}. Despite these advances, recent works such as B\"{o}hm et al. \cite{12} continue to model historic flood events statically, training models to recognize where water was, rather than establishing a parametrizable relationship to predict where water will be under differing dynamic rainfall conditions.

\subsection{Synthetic Aperture Radar Ground Truth}
Active microwave sensing, specifically C-band SAR from Sentinel-1, serves as the standard for flood label generation due to its cloud-penetrating capabilities. The specular reflection of radar pulses off smooth standing water yields a characteristic low-backscatter signal, enabling automated water delineation \cite{13,14}. This study utilizes rigorously post-processed Sentinel-1 SAR masks as the definitive ground truth.

\section{Proposed System Architecture}

\begin{figure}[htbp]
\centerline{\includegraphics[width=\linewidth]{images/system_architecture.png}}
\caption{The end-to-end GeoAI predictive framework, detailing data acquisition, geomorphic intelligence, the AI Core (Attention U-Net), and the presentation layer.}
\label{fig:architecture}
\vspace{-3mm}
\end{figure}

\subsection{Formal Problem Definition}
The mapping task is formalized as a dense pixel-wise binary segmentation problem. Let $\mathbf{X} \in \mathbb{R}^{B \times 7 \times H \times W}$ denote the multi-channel spatial input tensor, where $H=W=64$ represents the patch dimensions at 30m resolution. Let $\mathbf{Y} \in \{0,1\}^{B \times 1 \times H \times W}$ represent the binary SAR-derived flood label. The system learns a nonlinear mapping $f_\theta$ to predict pixel-wise flood susceptibility probabilities:
\begin{equation}
\hat{\mathbf{Y}} = f_\theta(\mathbf{X}), \quad \hat{\mathbf{Y}} \in [0,1]^{B \times 1 \times H \times W}
\end{equation}

\subsection{Feature Engineering and Tensor Construction}
All layers are resampled and aligned to a unified 30-meter EPSG:32643 coordinate reference grid.

\subsubsection{Terrain Morphology ($\mathbf{X}_0, \mathbf{X}_1, \mathbf{X}_2$)}
\textbf{Elevation (DEM)} from the SRTM 30m dataset ($\mathbf{X}_0$) provides hydrostatic potential.
\textbf{Slope} ($\mathbf{X}_1$) is computed via the Zevenbergen-Thorne discrete finite difference algorithm:
\begin{equation}
S = \arctan\left(\sqrt{\left(\frac{\partial z}{\partial x}\right)^2 + \left(\frac{\partial z}{\partial y}\right)^2}\right)
\end{equation}
\textbf{Topographic Wetness Index (TWI)} ($\mathbf{X}_2$) quantifies the spatial tendency to accumulate water, based on specific catchment area ($A_s$) and slope ($\beta$) \cite{15}:
\begin{equation}
\text{TWI} = \ln\left(\frac{A_s}{\tan \beta + \varepsilon}\right), \quad \varepsilon = 10^{-6}
\end{equation}

\subsubsection{Surrogate Drainage Hydraulics ($\mathbf{X}_3, \mathbf{X}_4$)}
To bypass the lack of municipal blueprints, a synthetic surrogate network is derived mathematically via the established D8 single-flow-direction algorithm \cite{16,21}.
From this network, \textbf{Drainage Proximity} ($\mathbf{X}_3$) is computed as the Euclidean distance transform $D_{\text{drain}}$ to the nearest channel pixel set $\mathcal{C}$.
\textbf{Drainage Density} ($\mathbf{X}_4$) relies on a uniform convolution representing the localized channel fraction. It is computed using a uniformly weighted $33 \times 33$ spatial kernel $K$ over the binary channel mask $\mathbf{C}$:
\begin{equation}
\mathbf{X}_4 = K * \mathbf{C}, \quad K_i = \frac{1}{33 \times 33}
\end{equation}

\subsubsection{Urban Impermeability ($\mathbf{X}_5$) and API Rainfall Modulation ($\mathbf{X}_6$)}
\textbf{Built-up Density} ($\mathbf{X}_5$) represents localized impermeable coverage, derived via applying the same $33 \times 33$ averaging kernel over the 10m ESA WorldCover urban mask.

The \textbf{External API Rainfall Layer} ($\mathbf{X}_6$) forms the parameterized meteorological constraint of the architecture. Instead of developing a redundant numerical weather prediction model, this layer fundamentally acts as an API ingestion and spatialization endpoint. It natively retrieves localized 24-hour meteorological forecasts $r_k$ from established third-party regional sensing APIs (e.g., OpenWeatherMap). Recognizing that urban waterlogging is dictated by spatially variable micro-bursts, these discrete external telemetry points are interpolated across the grid using Inverse Distance Weighting (IDW) to generate a localized precipitation surface $R(x,y)$, which is then normalized against a historical maximum threshold ($r_{\max} = 300\text{ mm}$):
\begin{equation}
\mathbf{X}_{6}(x,y) = \frac{1}{r_{\max}} \frac{\sum_{k} r_k / d_k(x,y)^p}{\sum_{k} 1 / d_k(x,y)^p}
\end{equation}
where $d_k(x,y)$ is the distance from pixel $(x,y)$ to external station coordinate $k$, and $p=2$ controls the decay rate. Rather than simulating complex continuous hydrodynamic rainfall-runoff interactions, this approach efficiently parameterizes the spatial intensity of the incoming storm load. It allows the deep network to learn a direct geographical correlation between variable micro-burst peak intensities and the geometric topology of eventual inundation, providing a computationally lightweight operational proxy for data-scarce regions.

\section{Attention U-Net Methodology}
An Attention U-Net architecture \cite{11} is employed to focus capacity on hydrologically critical terrain.

\subsection{Channel Attention and Spatial Attention}
A Squeeze-and-Excitation pattern implicitly weights the 7 varying input channels by mapping global characteristics through a Multi-Layer Perceptron (MLP):
\begin{equation}
CA(\mathbf{x}) = \sigma(\text{MLP}(\text{AvgPool}(\mathbf{x})) + \text{MLP}(\text{MaxPool}(\mathbf{x}))) \otimes \mathbf{x}
\end{equation}
During expansive decoding, skip-connection features $\mathbf{x}^l$ are modulated by a gating signal $\mathbf{g}$ to suppress irrelevant background regions (e.g., dry hillsides) and highlight relevant hydrologic features:
\begin{equation}
\alpha = \sigma(\text{ReLU}(\mathbf{W}_g \mathbf{g} + \mathbf{W}_x \mathbf{x}^l)), \quad \tilde{\mathbf{x}}^l = \alpha \otimes \mathbf{x}^l
\end{equation}

\subsection{Objective Function}
The network minimizes a pixel-wise Binary Cross-Entropy (BCE) loss:
\begin{equation}
\mathcal{L}_{\text{BCE}} = -\frac{1}{N} \sum_{i} \left[ y_i \log(\hat{y}_i) + (1 - y_i)\log(1 - \hat{y}_i) \right]
\end{equation}

\section{Experimental Setup}

\subsection{Dataset Configuration and Reproducibility Details}
The study utilizes rigorously post-processed 10m-resolution C-band Sentinel-1 SAR imagery acquired during the catastrophic August 2018 monsoon event (ascending pass), resampled to the 30m base grid. To resolve extreme class imbalance (flood pixels inherently constitute $<$15\% of landmass), a targeted patch-extraction strategy generated a dataset of exactly 2,000 spatial patches ($64 \times 64$ pixels, representing 3.68 $\text{km}^2$ per patch). The dataset was strictly balanced, containing exactly 1,000 flood-positive and 1,000 flood-negative (dry) geographical scenes to stabilize cross-entropy gradients. All models were trained and evaluated continuously on a single NVIDIA RTX 3060 (12GB VRAM) instance.

\subsection{Spatial Block Cross-Validation and Data Integrity}
Randomly shuffling image patches introduces a severe risk of data leakage due to spatial autocorrelation (Tobler's First Law of Geography). To explicitly prevent this, training datasets were generated using a \textbf{spatially disjoint block cross-validation} strategy. The 3,068 km$^2$ district was partitioned into non-overlapping geographic quadrants. Models were trained strictly on specific quadrants and evaluated on entirely disjoint regional blocks. To ensure convergence stability, patches within training blocks were balanced to maintain equal prior probabilities for flood/non-flood classes.

\subsection{Baseline Architectures and Hyperparameters}
The Attention U-Net ($\sim$7.81 million parameters, depth of 4 encoder-decoder blocks with standard 64-128-256-512 filter progression) was trained using an Adam optimizer ($\beta=(0.9, 0.999)$) with an initial learning rate of $0.001$, annealed via a Cosine Scheduler. Training executed over 30 epochs with a batch size of 16.

It is benchmarked against Traditional Machine Learning models (Logistic Regression, Support Vector Machine, and Random Forest evaluated on spatially flattened 7-dimensional vectors) and a state-of-the-art spatial CNN benchmark, \textbf{DeepLabV3+} \cite{18}, known for exceptional semantic segmentation via atrous convolutions.

\section{Results \& Evaluation}

\subsection{Quantitative Baseline Comparison}
Table \ref{tab:results} details benchmarked performance on the geographically independent spatial validation set (threshold = 0.5).

\begin{table}[htbp]
\caption{Comparative Benchmarks (Spatial Block Validation Set)}
\begin{center}
\resizebox{\linewidth}{!}{
\begin{tabular}{|l|c|c|c|c|}
\hline
\textbf{Model} & \textbf{Precision} & \textbf{Recall} & \textbf{F1-Score} & \textbf{AUC} \\
\hline
Logistic Regression & 0.000 & 0.000 & 0.000 & 0.798 \\
SVM (RBF Kernel) & 0.000 & 0.000 & 0.000 & 0.884 \\
Random Forest & 1.000 & 0.095 & 0.174 & \textbf{0.920} \\
DeepLabV3+ \cite{18}& 0.701 & 0.615 & 0.655 & 0.910 \\
Standard U-Net & 0.648 & 0.573 & 0.608 & 0.871 \\
\textbf{Attention U-Net} & \textbf{0.712} & \textbf{0.632} & \textbf{0.670} & 0.919 \\
\hline
\end{tabular}
}
\label{tab:results}
\end{center}
\end{table}

The Random Forest achieves a numerically high ROC-AUC (0.920), demonstrating a common artifact in imbalanced geospatial analysis: by aggressively predicting non-flooded states, it secures high true-negative rates that inflate the curve, despite severe under-prediction of actual flood mass (Recall = 0.095).
Conversely, the Attention U-Net demonstrates superior structural segmentation (F1 = 0.670), retaining spatial contiguity across complex terrain. It slightly outperforms the heavier DeepLabV3+ architecture (F1 = 0.655), validating the efficacy of the targeted spatial attention mechanism for resolving narrow fluvial networks.

\subsection{Temporal Generalization Evaluation}
To evaluate predictive capability beyond a static target mapping task, the optimal model was tested on a temporally disjoint, independent precipitation event (the subsequent August 2019 monsoon) using distinct SAR ground truth pairs. On this unseen temporal scenario, the Attention U-Net sustained an F1-score of 0.642. This structural retention across varying storm dynamics supports the hypothesis that the network has learned a generalized correlation between the parameterized $X_6$ rainfall surface and inundation outcomes, rather than simply overfitting a single historical landscape geometry.

\subsection{Ablation Study: Feature Significance}
An exhaustive leave-one-out feature ablation study was conducted. Removing the interpolated \textbf{Rainfall Channel} caused a 13.7\% reduction in F1-score. This confirms the network's utilization of the dynamic meteorological constraint as an active operational parameter. Removing the \textbf{Surrogate Drainage} feature yielded the most severe degradation (--20.7\% F1), empirically demonstrating that derivation of synthetic hydraulic pathways (D8 routing) effectively compensates for missing municipal blueprints within deep learning applications.

\subsection{Real-Time Operational Deployment}
The developed Hybrid CNN+Hydrological Physics model provides near real-time, dynamic flood susceptibility predictions based on varying meteorological inputs. To validate the interactive forecasting capabilities of the system, a scenario-based analysis was conducted for the Ernakulam district using the integrated web dashboard. The spatial distribution of flood probability was evaluated under three distinct rainfall intensity scenarios: baseline (0 mm), moderate (50 mm), and extreme (175 mm).

\subsubsection{Baseline Condition Analysis (0 mm)}
Under baseline conditions, representing a dry scenario with 0 mm rainfall intensity, the model correctly predicts negligible flood probability across the study area. The entire district exhibits very low susceptibility (visualized in green), confirming that the model has effectively learned to isolate rainfall as a primary dynamic trigger for inundation. This performance successfully avoids false-positive predictions in the absence of precipitation events, establishing a reliable ``zero-state'' for the system.

\begin{figure}[htbp]
\centerline{\includegraphics[width=\linewidth]{images/dashboard_5mm.png}}
\caption{Baseline Scenario (0 mm)}
\label{fig:dash_low}
\end{figure}

\subsubsection{Moderate Rainfall Response (50 mm)}
When the rainfall intensity is increased to 50 mm, representing a moderate precipitation event, the model demonstrates a physically consistent localized response. The flood probability increases in specific vulnerable zones (visualized in orange/red transitions), particularly along low-lying areas, river banks, and regions with high topographic wetness indices. This dynamic sensitivity highlights the successful integration of hydrological physics, allowing the CNN to adjust its base geomorphological susceptibility mapping according to precise rainfall inputs.

\begin{figure}[htbp]
\centerline{\includegraphics[width=\linewidth]{images/dashboard_56mm.png}}
\caption{Moderate Scenario (50 mm)}
\label{fig:dash_med}
\end{figure}

\subsubsection{Extreme Event Simulation (175 mm)}
Simulating an extreme rainfall event of 175 mm characteristic of the severe monsoonal downpours often experienced in Kerala reveals widespread and severe flood susceptibility (visualized in dark red/purple) across the majority of the Ernakulam district. The rapid expansion of high-probability zones correctly correlates with historical observations of massive inundation during extreme weather events. The model's ability to smoothly and logically scale the susceptibility output in direct proportion to the simulated rainfall intensity validates the robustness of the Hybrid CNN-Physics architecture for dynamic early warning applications.

\begin{figure}[htbp]
\centerline{\includegraphics[width=\linewidth]{images/dashboard_175mm.png}}
\caption{Extreme Scenario (175 mm)}
\label{fig:dash_high}
\end{figure}

\subsubsection{Conclusion on Interactive Utility}
Ultimately, the interactive scenario analysis confirms that the developed GeoAI system is highly responsive to meteorological variables. Unlike static flood risk maps, this approach provides stakeholders with a powerful tool to simulate custom weather events and visualize the corresponding localized impact. This functionality significantly enhances proactive disaster management and emergency response planning by offering a predictive look at potential crises before they manifest.

\section{Discussion \& Limitations}
The proposed framework addresses the critical data scarcity faced by developing nations by utilizing synthetic derivations and spatial interpolation. The successful temporal validation proves the API-driven approach moves beyond static hindcasting.

However, significant limitations remain. The D8 routing methodology models strictly surface-driven topological flow. Consequently, the framework remains inherently blind to subsurface storm-water pipe infrastructure. While acceptable for peri-urban or unmapped regions experiencing pluvial terrain flooding, its physical modeling capability is diminished in highly concretized metropolitan cores where underground pressurized hydraulics supersede surface gradients. Furthermore, treating time as discrete snapshots rather than a continuous sequence ignores hydraulic wave propagation speeds; future work must transition this spatial tensor network into a recurrent framework (e.g., ConvLSTM) to resolve sub-daily temporal dynamics.

\section{Conclusion}
This research presented a GeoAI framework engineered to provide actionable flood forecasting despite severe structural data deficits. By generating surrogate drainage networks via D8 topography and interpolating dynamic weather API telemetry into a 7-channel multi-modal tensor, the model achieves real-time interactive mapping capability. 

Evaluated under strict spatial cross-validation and temporally independent testing, the Attention U-Net demonstrated strong discriminative performance (F1-Score 0.670, ROC-AUC 0.919), remaining competitive with heavy agnostic algorithms while outperforming standard semantic architectures relative to spatial continuity. Feature ablation empirically validated the essential role of both the generated drainage proximity and the parameterized meteorological surface constraint. Serving via an optimized FastAPI backend, this open-source framework offers immediate, scalable assessment utility for regions lacking complex hydrodynamic infrastructure data.

\section*{Acknowledgement}
The authors would like to express their sincere gratitude to the management and Principal, Dr. Saju P John, of Providence College of Engineering, for providing the necessary infrastructure and a supportive environment to carry out this project. Appreciation is also extended to Dr. Prem Sankar C., Head of the Department of Artificial Intelligence and Project Coordinator, for his timely suggestions and overall coordination. The authors acknowledge all those who directly or indirectly contributed to the successful completion of this work.

\begin{thebibliography}{00}
\bibitem{1} D. Guha-Sapir, P. Hoyois, and R. Below, ``Annual Disaster Statistical Review 2015: The Numbers and Trends,'' \textit{CRED}, Brussels, 2016.
\bibitem{2} W. D. Shuster et al., ``Impacts of impervious surface on watershed hydrology: A review,'' \textit{Urban Water Journal}, vol. 2, pp. 263--275, 2005.
\bibitem{3} P. Lal et al., ``Evaluating the 2018 extreme flood hazard events in Kerala, India,'' \textit{Remote Sensing Letters}, vol. 11, no. 5, pp. 436--445, 2020.
\bibitem{4} V. A. Rangari et al., ``Assessment of inundation risk in urban floods using HEC RAS 2D,'' \textit{Modeling Earth Systems and Environment}, vol. 5, no. 4, pp. 1839--1851, 2019.
\bibitem{5} M. S. Tehrany et al., ``Flood susceptibility mapping using novel ensemble models in GIS,'' \textit{Journal of Hydrology}, vol. 512, pp. 332--343, 2014.
\bibitem{6} K. Khosravi et al., ``A GIS-based flood susceptibility assessment,'' \textit{Natural Hazards}, vol. 83, pp. 947--987, 2016.
\bibitem{7} G. Konapala et al., ``Exploring Sentinel-1 and Sentinel-2 diversity for flood inundation mapping using deep learning,'' \textit{ISPRS Journal}, vol. 180, pp. 163--173, 2021.
\bibitem{8} S. Kabir et al., ``A deep convolutional neural network model for rapid prediction of fluvial flood inundation,'' \textit{Journal of Hydrology}, vol. 590, p. 125481, 2020.
\bibitem{9} W. Zhao et al., ``GeoAI driven spatial analysis of extreme floods,'' \textit{IEEE Trans. on Geoscience and Remote Sensing}, vol. 60, pp. 1-14, 2022.
\bibitem{10} N. Kazakis et al., ``Assessment of flood hazard areas at a regional scale,'' \textit{Science of the Total Environment}, vol. 538, pp. 555--563, 2015.
\bibitem{11} O. Oktay et al., ``Attention U-Net: Learning Where to Look for the Pancreas,'' \textit{Medical Imaging with Deep Learning}, 2018.
\bibitem{12} L. B\"{o}hm et al., ``FloodTransNet: Transformer-based flood extent mapping,'' \textit{Int. J. Applied Earth Observation}, vol. 119, 2023.
\bibitem{13} S. Schlaffer et al., ``Flood detection from multi-temporal SAR data,'' \textit{Int. J. Applied Earth Observation}, vol. 38, pp. 15--24, 2015.
\bibitem{14} P. Manjusree et al., ``Optimization of threshold ranges for rapid flood mapping,'' \textit{Int. J. Disaster Risk Science}, vol. 3, pp. 113--122, 2012.
\bibitem{15} G. J. Beven and M. J. Kirkby, ``A physically based, variable contributing area model of basin hydrology,'' \textit{Hydrological Sciences Bulletin}, vol. 24, pp. 43--69, 1979.
\bibitem{16} B. Pradhan and M. I. Youssef, ``A 100-year maximum flood susceptibility mapping,'' \textit{Journal of Flood Risk Management}, vol. 4, pp. 189--202, 2011.
\bibitem{17} P. Jiang et al., ``Spatial analysis and deep learning with Swin-Unet for flood mapping,'' \textit{Water Resources Research}, vol. 58, 2022.
\bibitem{18} L.-C. Chen et al., ``Encoder-Decoder with Atrous Separable Convolution for Semantic Image Segmentation,'' \textit{ECCV}, 2018.
\bibitem{19} X. Shi et al., ``Convolutional LSTM Network: A Machine Learning Approach for Precipitation Nowcasting,'' \textit{NeurIPS}, 2015.
\bibitem{20} T. Kipf and M. Welling, ``Semi-Supervised Classification with Graph Convolutional Networks,'' \textit{ICLR}, 2017.
\bibitem{21} J. F. O'Callaghan and D. M. Mark, ``The extraction of drainage networks from digital elevation data,'' \textit{Computer Vision, Graphics, and Image Processing}, vol. 28, no. 3, pp. 323-344, 1984.
\end{thebibliography}

\end{document}
